\chapter{Massive Weak Exchange, Virtual Particles, and Broken Symmetry}

The Standard Model group becomes much less mysterious when we read it as a
set of allowed motions on the fermion table. Color transformations move
vertically among red, green, and blue. Charged weak transformations move
horizontally within quark and lepton doublets. Electromagnetic phase
transformations supply a third operation. The corresponding gauge bosons are
not decorations added after the symmetries have been found: emitting a gauge
boson physically performs the same change of state that a generator of the
symmetry describes.

That observation sets up the main question of this lecture. A weak vertex has
a coupling comparable in size to an electromagnetic vertex, yet beta decay is
enormously slower than a typical strong or electromagnetic process. The
missing ingredient is propagation. At low momentum, exchanging a massive
\(W\) costs a factor of order \(1/m_W^2\) in the amplitude. We shall develop
that statement first in distance space and then in momentum space, follow its
consequences through a long discussion of real and virtual particles, and
finally turn to the first qualitative model of spontaneous symmetry breaking.

\section{Reading Gauge Symmetry from the Fermion Table}

\subsection{Color acts vertically}

Arrange the six quark flavors in three weak families,
\begin{equation}
  (u,d),\qquad(c,s),\qquad(t,b),
\end{equation}
and give every quark a red, green, or blue color label. The resulting quark
table has three color rows and six flavor columns. Color \(SU(3)\) mixes the
three entries in any one flavor column. For example, one generator may turn a
red up quark into a green up quark. In particle language the same operation is
accompanied by emission of a gluon carrying the compensating color and
anticolor.

The same color transformation must act on every flavor. The gluon that
changes a red up quark into a green up quark is the same species of gluon that
changes a red top quark into a green top quark. A genuine symmetry operation
does not apply to the first generation while ignoring the others. Thus color
motion is vertical and universal:
\begin{equation}
  q_r\longleftrightarrow q_g\longleftrightarrow q_b,
  \qquad q\in\{u,d,c,s,t,b\}.
\end{equation}
It changes color but not flavor.

This is the characteristic relation between a gauge group and its gauge
bosons. The infinitesimal generators describe transformations among states,
and emission or absorption of the corresponding bosons realizes those
transformations as physical processes.

\subsection{Weak transformations act horizontally}

The charged weak interaction moves in the other direction. An up-type quark
can emit a charged weak boson and become its down-type partner without
changing color:
\begin{equation}
  u_b\longrightarrow d_b+W^+.
\end{equation}
Exactly the same \(W^+\) occurs in
\begin{equation}
  c_g\longrightarrow s_g+W^+,
  \qquad
  t_r\longrightarrow b_r+W^+.
\end{equation}
The labels \(r,g,b\) are spectators. Gluons do not move horizontally across
the table, and charged \(W\) bosons do not move vertically through color.

Add the electromagnetic phase symmetry \(U(1)\), under which every charged
field acquires a phase. The first-pass group structure is then
\begin{equation}
  SU(3)\times SU(2)\times U(1).
\end{equation}
The cross signs mean that these are independent factors: one acts on color,
one on weak doublets, and one by a phase. This is the useful rough picture,
but it is not yet the complete electroweak story.\editorialnote{The exact Standard Model gauge group is conventionally written
\(SU(3)_c\times SU(2)_L\times U(1)_Y\). The weak \(SU(2)_L\) acts only on
left-handed doublets, and the final electromagnetic \(U(1)_{\rm em}\) emerges
from \(SU(2)_L\times U(1)_Y\) after symmetry breaking. The lecture explicitly
warns that its first-pass identification of the last two factors is not yet
fully consistent; later lectures supply the missing chiral and electroweak
structure.}

\subsection{The same pattern includes leptons}

Below the three color rows place a colorless lepton row:
\begin{equation}
\begin{array}{c|cc|cc|cc}
 &u&d&c&s&t&b\\ \hline
Q&+\frac23&-\frac13&+\frac23&-\frac13&+\frac23&-\frac13\\
\text{lepton partner}
 &\nu_e&e^-&\nu_\mu&\mu^-&\nu_\tau&\tau^-\\
Q&0&-1&0&-1&0&-1
\end{array}
\end{equation}
Within every adjacent pair the electric charges differ by one unit. That is
why the same charged weak boson can connect \(u\) to \(d\), \(c\) to \(s\),
\(t\) to \(b\), and a charged lepton to its neutrino. The \(W\) vertex cares
about the appropriate quantum numbers and charge difference, not about
whether the fermion is a quark or a lepton.

\begin{figure}[htbp]
\centering
\includegraphics[width=0.88\textwidth]{lecture_06_figure_07.png}
\caption{The completed fermion table: three quark-color rows above the
colorless lepton row, with the quark charges and three lepton doublets visible.}
\lectureframe{00:18:55}
\end{figure}

The muon has the same charge and known intrinsic quantum numbers as the
electron, but is about two hundred times heavier. The tau is heavier still.
The neutrino masses are very small and were not known in detail in the
lecture's discussion. No principle has yet emerged here that explains the
observed pattern of family masses.

\begin{classroomqa}
\audiencequestion{Why are leptons, baryons, and mesons called by those names?
\lecturetimestamp{00:19:02}}
\lecturerresponse{The historical Greek roots meant roughly light, heavy, and
intermediate. This mass-based naming became misleading. The muon was once
called the mu meson because its mass lay between the electron and proton
masses, but its interactions show that it is a lepton, much more like an
electron than a pion. Today a meson means a quark--antiquark hadron and is a
boson; the muon is not a meson.}
\end{classroomqa}

\section{The Puzzle of Weakness}

\subsection{A concrete charged-current process}

Use neutron beta decay as the working diagram. The valence quark
compositions are
\begin{equation}
  n=udd,\qquad p=uud.
\end{equation}
One down quark makes the transition
\begin{equation}
  d\longrightarrow u+W^-,
\end{equation}
and the charged boson produces
\begin{equation}
  W^-\longrightarrow e^-+\bar\nu_e.
\end{equation}
Together,
\begin{equation}
  n\longrightarrow p+e^-+\bar\nu_e.
\end{equation}
The other two valence quarks carry through as spectators. At the first vertex
the neutron becomes a proton; at the second, the same weak interaction that
connects a lepton doublet converts the \(W^-\) into the observed electron and
electron antineutrino.

\subsection{Three couplings, three very different impressions}

In QED the electron charge \(e\) is the vertex amplitude. Squaring it gives
the characteristic emission probability, with conventional factors collected
in
\begin{equation}
  \alpha=\frac{e^2}{4\pi}\simeq\frac1{137}
\end{equation}
in natural units. The strong interaction has a coupling \(g_s\), and at a
typical hadronic scale its analogue \(\alpha_s\) may be of order \(0.2\).
That is much larger than \(1/137\), so gluon emission is common when a quark
is struck.

Let \(g_w\) denote the charged weak coupling. Its square is not fantastically
small; at the order-of-magnitude level used here it is comparable to the
electromagnetic coupling. Yet a free neutron waits roughly fifteen minutes
before decaying, while a strong decay may occur on a hadronic crossing time.
The historical estimate in the lecture uses the older rough figure of twelve
minutes, but either number makes the same puzzle unavoidable.\editorialnote{The modern free-neutron mean lifetime is about
\(8.8\times10^2\,{\rm s}\), with ongoing experimental tension between major
measurement methods. The lecture's rough lifetime is retained as the
pedagogical scale, not as a precision value.}
\begin{quote}
If \(g_w\) is not extraordinarily small, what makes weak processes weak?
\end{quote}

The answer cannot lie at the vertices alone. Every Feynman diagram also
contains propagators, the factors that describe the intermediate motion from
one vertex to another.

\section{A Massive Propagator in Distance Space}

\subsection{Vertices are not the whole diagram}

The beta-decay diagram contains two weak vertices, contributing two powers of
\(g_w\), and a \(W\) line joining them. Denote the position-space propagator
by \(\Delta(x-y)\), where \(x\) is the emission event and \(y\) the absorption
or conversion event. It is the amplitude kernel associated with carrying the
field disturbance between those events.

The exact relativistic propagator contains more structure than we need. Its
important qualitative feature is its dependence on mass. At short spacelike
distance it has a power-law singularity, schematically
\begin{equation}
  \Delta(d)\sim\frac1{d^2}.
\end{equation}
A massless field retains a power-law tail. A massive field, by contrast, is
exponentially suppressed once the separation becomes large compared with its
Compton scale. The graph therefore bends sharply downward.

\subsection{Natural units locate the bend}

Set
\begin{equation}
  c=\hbar=1.
\end{equation}
Then
\begin{equation}
  [m]=[E]=[p]=L^{-1}.
\end{equation}
The familiar relations explain each equality: \(E=mc^2\) identifies mass and
energy, while \(p\,x\) has the dimensions of \(\hbar\), so momentum and
distance have inverse units. The only length that can be formed from the mass
\(m\) is therefore
\begin{equation}
  d_m\sim\frac1m.
\end{equation}
Restoring constants gives
\begin{equation}
  d_m\sim\frac{\hbar}{mc},
\end{equation}
the Compton wavelength. A heavier exchanged particle has a shorter effective
range.

\begin{figure}[htbp]
\centering
\includegraphics[width=0.78\textwidth]{lecture_06_figure_02.png}
\caption{Natural-unit dimensions and the qualitative propagator falloff. The
black curve continues with a power-law tail; the red massive curve turns down
near the inverse-mass scale.}
\lectureframe{00:27:19}
\end{figure}

\begin{figure}[htbp]
\centering
\begin{tikzpicture}[x=1.0cm,y=0.9cm]
  \draw[->] (0,0) -- (7,0) node[right] {$d$};
  \draw[->] (0,0) -- (0,5) node[above] {$|\Delta(d)|$};
  \draw[thick]
    plot[smooth] coordinates {(0.45,4.5) (0.9,3.3) (1.6,2.5)
      (2.5,1.9) (3.7,1.45) (5.1,1.12) (6.4,0.9)};
  \draw[thick,red]
    plot[smooth] coordinates {(0.45,4.5) (0.9,3.3) (1.6,2.5)
      (2.5,1.9) (3.2,1.55) (3.7,0.95) (4.2,0.36) (4.8,0.1)};
  \draw[dashed] (3.4,0) -- (3.4,1.4);
  \node at (5.7,1.4) {$m=0$};
  \node[red] at (4.75,0.72) {$m>0$};
  \node[below] at (3.4,0) {$d\sim m^{-1}$};
\end{tikzpicture}
\caption{A clean schematic of the board argument. It is not an exact formula
for a four-dimensional propagator; it isolates the mass-controlled range.}
\end{figure}

\begin{classroomqa}
\audiencequestion{What is plotted on the vertical axis?
\lecturetimestamp{00:28:30}}
\lecturerresponse{The size of the propagator, denoted by \(\Delta\). Its
square contributes to a probability after the complete amplitude, external
states, and phase-space factors are included. The graph is meant to show that
the amplitude for massive exchange over a long distance shuts off rapidly.\editorialnote{A propagator by itself is not a directly measurable
probability amplitude for a particle following a classical path. The lecture's
language is a useful range estimate: the propagator is one factor in the full
scattering or decay amplitude.}}
\end{classroomqa}

\section{Fourier Transform and the Fermi Limit}

\subsection{Experiments specify momenta}

An experiment usually measures external momenta rather than the two
spacetime points at which an internal line begins and ends. In beta decay, the
momentum carried by the \(W\) is inferred from the difference between the
neutron and proton momenta, and it equals the total momentum of the electron
and antineutrino. Position-space and momentum-space propagators are related by
a Fourier transform.
\begin{equation}
  k=p_n-p_p=p_e+p_{\bar\nu_e}.
\end{equation}

In the positive-denominator convention used on the board,
\begin{equation}
  \widetilde{\Delta}_W(k)=\frac1{k^2+m_W^2},
  \qquad
  \widetilde{\Delta}_\gamma(k)=\frac1{k^2}.
\end{equation}
The massless photon propagator becomes large at small momentum transfer. The
massive \(W\) propagator never exceeds a scale of order \(1/m_W^2\). For beta
decay,
\begin{equation}
  |k^2|\ll m_W^2,
\end{equation}
so the details of \(k^2\) are negligible and
\begin{equation}
  \widetilde{\Delta}_W(k)\simeq\frac1{m_W^2}.
\end{equation}

\begin{classroomqa}
\audiencequestion{Should the relativistic \(k^2\) contain a minus sign between
energy and spatial momentum?
\lecturetimestamp{00:35:11}}
\lecturerresponse{Yes. With a mostly-minus Minkowski metric,
\(k^2=(k^0)^2-\boldsymbol{k}^2\), and the scalar denominator is conventionally
\(k^2-m_W^2+i\epsilon\). The board is using the equivalent positive
spacelike or Euclidean magnitude. In low-energy beta decay the transferred
four-momentum is tiny compared with \(m_W\), so either notation gives the same
\(1/m_W^2\) suppression up to an irrelevant sign in the amplitude.}
\end{classroomqa}

\subsection{The low-energy amplitude}

Two weak vertices and one propagator give
\begin{align}
  \mathcal A_{\rm weak}
  &\sim g_w\,\widetilde{\Delta}_W(k)\,g_w\\
  &\sim\frac{g_w^2}{k^2+m_W^2}
  \simeq\frac{g_w^2}{m_W^2}.
\end{align}
Squaring the amplitude produces the characteristic suppression
\begin{equation}
  |\mathcal A_{\rm weak}|^2
  \sim\frac{g_w^4}{m_W^4},
\end{equation}
before the kinematic and phase-space factors appropriate to the particular
decay are included. The weak interaction is weak at nuclear energies not
because its elementary gauge coupling is absurdly small, but because its
mediator is heavy.

\begin{figure}[htbp]
\centering
\includegraphics[width=0.86\textwidth]{lecture_06_figure_03.png}
\caption{The neutron--proton weak process and the momentum transfer carried
by the internal line, as discussed while fixing the propagator sign.}
\lectureframe{00:35:40}
\end{figure}

\begin{figure}[htbp]
\centering
\includegraphics[width=0.84\textwidth]{lecture_06_figure_04.png}
\caption{The beta-decay topology beside the boxed low-energy estimate
\(g_w^2/m_W^2\).}
\lectureframe{00:37:25}
\end{figure}

Numerically \(m_W\) is roughly one hundred proton masses at the accuracy of
this lecture. A factor of \(m_W^{-2}\) in the amplitude becomes
\(m_W^{-4}\) in a rate, so even a coupling of ordinary gauge-theory size
produces a very long lifetime.

\begin{classroomqa}
\audiencequestion{Are we talking about a suppression of one in a million or
one in a billion, and how was the \(W\) mass originally known?
\lecturetimestamp{00:38:11}}
\lecturerresponse{The lecture's rough counting points toward the billion-scale
smallness of a low-energy weak probability relative to an unsuppressed
hadronic process. Historically, low-energy experiments did not separately
measure \(g_w\) and \(m_W\). They measured the small combination
\(g_w^2/m_W^2\), called the Fermi constant. Once the gauge coupling was tied
to the electroweak theory, that combination predicted a very heavy mediator;
later accelerators produced the \(W\) directly and measured its mass.}
\end{classroomqa}

At momenta far below \(m_W\), every experiment sees nearly the same constant:
\begin{equation}
  G_F\sim\frac{g_w^2}{m_W^2}.
\end{equation}
This is why Fermi could describe beta decay in the 1930s as a pointlike
four-fermion interaction, long before the mediator was observed. As collision
energies increased, experiments became sensitive to the momentum dependence
hidden in the propagator, confirming that the effective constant was the
low-energy shadow of massive exchange.\editorialnote{With standard electroweak normalization at tree level,
\[
  \frac{G_F}{\sqrt2}=\frac{g^2}{8m_W^2}.
\]
The lecture deliberately suppresses numerical and spinor factors and keeps
only the scaling \(G_F\sim g_w^2/m_W^2\). The \(W\) and \(Z\) were directly
discovered at CERN in 1983; the lecture's passing mid-to-late-1970s dating is
too early for the direct observation.}

Historically, the idea that weak interactions might be mediated by a heavy
particle emerged in the mid-twentieth century; the lecture mentions Julian
Schwinger among the theorists who explored that possibility and notes that the
letter \(W\) probably stood for weak. Low-energy data already implied that
such a mediator would have to be extremely heavy on the nuclear scale, long
before accelerators had enough energy to make it as a real particle.

\subsection{Three independent controls on a rate}

A classroom question about nuclear binding separates three logically
different controls on a transition:
\begin{enumerate}
  \item vertex coupling constants;
  \item propagators, including masses and momentum dependence;
  \item available initial-to-final energy and phase space.
\end{enumerate}
Binding can close a decay channel because the final state would have greater
energy than the initial state. That is a kinematic prohibition, not a further
reduction of \(g_w\) or a change in \(m_W\). When a channel remains open but
has little available energy, its phase space can still be strongly
suppressed.

\begin{classroomqa}
\audiencequestion{Does the field or binding of a nearby proton suppress beta
decay through the same propagator mechanism?
\lecturetimestamp{00:41:13}}
\lecturerresponse{Not in the example being discussed. Binding matters because
it changes the energies of the initial and final nuclear states. A process
cannot occur if the final state is energetically inaccessible. Couplings,
propagators, and available energy are three separate ingredients.}
\end{classroomqa}

\begin{classroomqa}
\audiencequestion{Are the coupling constants really constant?
\lecturetimestamp{00:42:55}}
\lecturerresponse{Only at a fixed scale. Quantum corrections make effective
couplings run with momentum or distance. The present argument treats each
coupling at the scale relevant to the process; scale dependence will be
developed later.}
\end{classroomqa}

\section{Virtual Exchange without Abandoning Conservation Laws}

\subsection{Why the heavy line can appear}

Ordinary neutron decay does not contain enough energy to produce a free,
on-shell \(W\). The \(W\) in the diagram is internal and off shell. The
lecture introduces this fact with the traditional energy--time estimate
\begin{equation}
  \Delta E\,\Delta t\lesssim\hbar:
\end{equation}
an intermediate configuration whose classical energy appears wrong by
\(\Delta E\) can contribute only over a time of order
\(\hbar/\Delta E\).

The tunneling analogy makes the intended picture vivid. Put an alpha particle
in a nuclear potential well. Its energy lies below the top of the barrier, so
classically it cannot escape. Quantum mechanically the wavefunction extends
through the barrier, and the particle can emerge at a point on the far side
with the same total energy it had initially. If we try to describe the
interior of the barrier as an ordinary classical trajectory, the potential
energy seems too large and the kinetic energy would have to be negative.
That failed classical picture is analogous to trying to regard the internal
\(W\) as a real particle.

\begin{figure}[htbp]
\centering
\includegraphics[width=0.78\textwidth]{lecture_06_figure_06.png}
\caption{The tunneling analogy: the initial point and the marked point on the
far slope have equal total energy, while the intervening barrier is
classically forbidden.}
\lectureframe{00:46:00}
\end{figure}

\begin{figure}[htbp]
\centering
\begin{tikzpicture}[x=1.0cm,y=0.9cm]
  \draw[->] (0,0) -- (7,0) node[right] {$x$};
  \draw[->] (0,0) -- (0,4.8) node[above] {$V(x)$};
  \draw[thick]
    plot[smooth] coordinates {(0.3,1.4) (1.0,0.7) (1.7,0.55)
      (2.1,0.9) (3.2,3.9) (4.2,2.1) (5.0,0.75) (6.6,0.18)};
  \draw[dashed] (1.7,0.75) -- (5.0,0.75);
  \fill (1.7,0.75) circle (0.06);
  \fill (5.0,0.75) circle (0.06);
  \node at (3.35,0.35) {$\Delta E\,\Delta t\lesssim\hbar$};
\end{tikzpicture}
\caption{A readable reconstruction of the barrier and equal-energy points.}
\end{figure}

In the beta-decay picture the neutron may fluctuate into a proton plus an
internal \(W^-\), but the heavy line must be removed quickly. It may be
reabsorbed, restoring the neutron, or it may terminate at the second vertex
and produce the much lighter electron and antineutrino. The full initial and
final energies agree.\editorialnote{Energy and momentum are conserved at every Feynman vertex.
Virtual particles do not literally borrow energy and repay it later. The
precise statement is that an internal line is integrated over four-momenta
that need not satisfy the real-particle mass-shell relation
\(k^2=m_W^2\). The uncertainty-principle and tunneling language is retained
because it is the lecture's explanatory route, but the off-shell statement is
the modern QFT formulation.}

\subsection{Real W Bosons}

A sufficiently energetic accelerator can produce a real \(W\). It then has
its own proper lifetime, unrelated to a putative uncertainty time for the
virtual line. A boosted real \(W\) lives longer in the laboratory by ordinary
relativistic time dilation. In a two-body leptonic decay of a \(W\) at rest,
\begin{equation}
  W^-\longrightarrow e^-+\bar\nu_e,
\end{equation}
the electron and antineutrino carry large, nearly back-to-back kinetic
energies, each of order \(m_W/2\). The rest energy has not disappeared; it
has become final-state energy. Muon and tau channels, as well as quark
channels, are also available.

\begin{classroomqa}
\audiencequestion{How long does a real \(W\) live, and can its lifetime be
computed from the same coupling?
\lecturetimestamp{00:49:49}}
\lecturerresponse{Yes. One calculates each allowed decay amplitude, squares
it, integrates over final-state phase space, and sums the partial widths. The
inverse total width gives the lifetime. The transcript's spoken numerical
estimate is garbled, but the qualitative point is that the proper lifetime is
extremely short and a moving \(W\) is time-dilated in the usual way.\editorialnote{Using the measured width
\(\Gamma_W\simeq2.1\,{\rm GeV}\) gives
\(\tau_W=\hbar/\Gamma_W\simeq3.2\times10^{-25}\,{\rm s}\). This modern
number resolves the corrupted numerical exchange without attributing it to
the lecture.}}
\end{classroomqa}

\begin{classroomqa}
\audiencequestion{How does a squared amplitude become a rate and hence a
lifetime?
\lecturetimestamp{00:50:52}}
\lecturerresponse{The squared matrix element is combined with the density of
available final states, conservation delta functions, and normalization
factors to obtain a probability per unit time, the decay width \(\Gamma\).
The survival probability is \(e^{-\Gamma t}\) in units with \(\hbar=1\), so
\(\tau=1/\Gamma\), or \(\tau=\hbar/\Gamma\) with units restored. A small weak
width therefore means a long weak lifetime.}
\end{classroomqa}

\begin{classroomqa}
\audiencequestion{Can one assign a velocity to a particle tunneling across
the barrier?
\lecturetimestamp{00:52:42}}
\lecturerresponse{One can divide an inferred distance by an inferred transit
time, but that quantity should not be interpreted as the velocity of a
classical particle or a signal. Attempting to measure a trajectory inside the
barrier changes the state and destroys the original tunneling experiment.}
\end{classroomqa}

\section{What Does It Mean to Observe a Virtual Particle?}

\subsection{External particles and internal lines}

Real particles are the incoming and outgoing asymptotic states that can be
prepared or detected far from the interaction region. Virtual particles are
internal lines in a perturbative representation of the amplitude. They are
not additional hidden detector tracks.

\begin{classroomqa}
\audiencequestion{Is there a clear definition of the difference between a
real and a virtual particle?
\lecturetimestamp{00:55:05}}
\lecturerresponse{A real particle is an external on-shell state entering or
leaving the interaction region. A virtual particle is an internal off-shell
line in a Feynman diagram. The latter is part of the mathematical expansion,
not an independently observed object.}
\end{classroomqa}

To resolve a process lasting only \(\Delta t\), a probe needs frequencies of
order \(1/\Delta t\) and hence energy of order \(\hbar/\Delta t\). A photon
sharp enough to resolve an alpha particle beneath a high barrier can supply
the energy needed to lift it into the barrier region. Likewise, a probe
energetic enough to resolve the short-lived internal \(W\) can create a real
\(W\). The attempted observation turns the original process into a different
experiment.

The lecture phrases this as a recurring quantum-mechanical frustration: the
intervention required to check the forbidden intermediate picture supplies
the very energy that makes the apparently forbidden configuration real.

\subsection{Fluctuations of a proton}

In perturbation theory a proton state contains amplitudes for configurations
such as
\begin{equation}
  p\longleftrightarrow n+W^+,
\end{equation}
and a neutron similarly has
\begin{equation}
  n\longleftrightarrow p+W^-.
\end{equation}
If no other interaction intervenes, the internal boson is reabsorbed. A hard
probe can instead leave a \(W\) among the outgoing particles. One may call
that act knocking out a fluctuation or creating a \(W\). The language differs;
the measurable statement is the same final-state amplitude.

\begin{classroomqa}
\audiencequestion{When a hard probe produces an outgoing \(W\), did it knock
out a particle that was already there or create a new one?
\lecturetimestamp{01:01:20}}
\lecturerresponse{Either picture can organize the mathematics. One
time-ordered description says that the proton fluctuated to \(n+W^+\) and the
probe removed the neutron before reabsorption; another says the collision
created the final \(W\). What is invariant is the calculated probability for
the observed final state, not a unique classical story about what existed
before measurement.}
\end{classroomqa}

Increasing the collision energy does not expose a fixed inventory of objects
that had been hidden inside the target. It opens many possible final states:
one \(W\), several bosons, quark jets, hadrons, and other particles allowed by
the conserved quantum numbers. Striking even an elementary field hard enough
can produce protons and many other composite objects. What comes out need not
have been present beforehand as a constituent. The lecture pushes the point
to comic extremes: if energy and quantum numbers allow it, even tables,
chairs, or a hydrogen atom are not logically excluded as fantastically
unlikely final arrangements. The joke carries a serious warning against
defining elementarity by a naive inventory of collision debris.

\subsection{Vacuum fluctuations and temperature}

A fluctuation attached to a proton is not the same as a vacuum diagram. The
first is a component of an interacting proton state; a vacuum fluctuation has
no external proton line. In either case, a sufficiently energetic probe may
produce real particles corresponding to an allowed final state. Two energetic
photons, for example, can produce an electron--positron pair when the
center-of-mass energy exceeds threshold.

\begin{classroomqa}
\audiencequestion{Are these reactions dependent on temperature?
\lecturetimestamp{01:05:17}}
\lecturerresponse{The elementary amplitude does not require a thermal bath.
Temperature matters statistically: a high-temperature system contains many
particles with high kinetic energy, so threshold-crossing collisions become
common. Heating a system containing a bound alpha particle may also supply
enough energy for over-the-barrier escape in addition to quantum tunneling.
A few isolated energetic particles constitute a beam; many thermally
distributed energetic particles constitute a hot bath.}
\end{classroomqa}

\begin{classroomqa}
\audiencequestion{Is the \(W\) emitted in a proton fluctuation a fluctuation
of the vacuum?
\lecturetimestamp{01:06:56}}
\lecturerresponse{No. The internal \(W\) in
\(p\leftrightarrow n+W^+\) is a fluctuation of the interacting proton state.
A closed virtual process with no external proton is a vacuum fluctuation.
They obey the same field-theory rules but belong to different amplitudes.}
\end{classroomqa}

\begin{classroomqa}
\audiencequestion{Why is the antineutrino required at the leptonic \(W\)
vertex?
\lecturetimestamp{01:08:51}}
\lecturerresponse{It supplies the conserved lepton quantum numbers and is part
of the basic charged-current vertex connecting \(e\), \(\nu_e\), and \(W\).
The bar distinguishes the outgoing electron antineutrino from a neutrino.}
\end{classroomqa}

\section{How Far Can Energy Conversion Be Pushed?}

The class next drives the rule that energy can create particles to an
intentionally absurd limit. Could a coffee bean become a chocolate cake?
There is no practical process of that kind, but relativistic quantum theory
does not impose a rule that final matter must resemble the initial object. It
requires enough energy, the correct conserved charges, and a nonzero
amplitude. A kilogram has an energy of order
\begin{equation}
  Mc^2\sim9\times10^{16}\ {\rm J},
\end{equation}
so any microscopic collision that assembled a particular kilogram-scale
object would be beyond fantastically improbable.

The Planck mass,
\begin{equation}
  m_{\rm P}=\sqrt{\frac{\hbar c}{G}}
  \simeq2.2\times10^{-5}\ {\rm g},
\end{equation}
is not a maximum possible mass. It is the scale at which gravity becomes
strong for a quantum localized to its Compton wavelength. A kilogram contains
roughly \(5\times10^7\) Planck masses.

\begin{classroomqa}
\audiencequestion{If collisions exceed the Planck scale, can arbitrarily
heavy objects such as chairs or planets emerge?
\lecturetimestamp{01:10:37}}
\lecturerresponse{At trans-Planckian center-of-mass energy concentrated in a
small region, gravitational collapse is the overwhelmingly natural outcome:
a black hole. A highly organized macroscopic object would be an
extraordinarily rare fluctuation. The lecture invokes time reversal: because
a chair can fall into a black hole, the reverse amplitude is not exactly
zero. Its probability is suppressed by an enormous entropy factor, so
possibility must not be confused with practical likelihood.}
\end{classroomqa}

This excursion reinforces the earlier warning. Quantum field theory predicts
amplitudes between specified initial and final states. A picturesque claim
that every product was literally contained in the target before impact is
not part of the prediction.

\section{Bosons, Their Interactions, and Their Statistics}

\subsection{Photons can scatter}

Being a boson does not mean being noninteracting. Boson or fermion specifies
exchange statistics; interaction vertices are separate dynamical facts.
Classical light beams normally pass through one another because QED has no
tree-level four-photon vertex and low-energy photon--photon scattering is
extremely small. Nevertheless, two photons can interact through a virtual
electron--positron loop:
\begin{equation}
  \gamma+\gamma
  \longrightarrow(e^-e^+)_{\rm virtual}
  \longrightarrow\gamma+\gamma.
\end{equation}
When the center-of-mass energy is high enough, they may instead create a real
pair,
\begin{equation}
  \gamma+\gamma\longrightarrow e^-+e^+.
\end{equation}
Optical photons have energies of a few electron volts, while the pair
threshold is of order \(2m_e c^2\approx1.02\,{\rm MeV}\). The virtual process
is therefore heavily suppressed at ordinary optical energies, which is why
crossing light beams appear not to notice each other.

\begin{classroomqa}
\audiencequestion{Do virtual \(W\) bosons interact with other bosons simply
because they are bosons?
\lecturetimestamp{01:12:47}}
\lecturerresponse{They can interact, but not merely because they obey Bose
statistics. Gauge vertices determine the dynamical interactions. Statistics
instead determines how multiparticle amplitudes behave when identical
particles are exchanged. Photons interact weakly through charged virtual
particles, and the electroweak theory also contains direct non-Abelian
\(W\)-boson self-interactions.}
\end{classroomqa}

\subsection{Bose enhancement is not a force}

Fermion versus boson statistics becomes conspicuous when several identical
particles occupy nearby quantum states. Bosons favor multiple occupation;
fermions exclude identical occupation. In a typical high-energy collision,
particles fly into well-separated momentum states, so this statistical
effect may be less obvious than it is in a laser or condensate.

In principle, \(W\) bosons possess the same Bose enhancement. The lecture
jokes that one could imagine a system that lases \(W\) bosons. Their enormous
mass and ultrashort lifetime make such a device fantastically impractical,
but the statistics is indeed bosonic.

\begin{classroomqa}
\audiencequestion{What color is a proton or neutron?
\lecturetimestamp{01:18:41}}
\lecturerresponse{A physical hadron is a color singlet, often called
colorless or white. The shorthand says that a baryon contains one red, one
green, and one blue quark. More precisely, its color wavefunction is the
antisymmetric superposition
\[
  |{\rm singlet}\rangle
  =\frac1{\sqrt6}\epsilon_{abc}|q^a q^b q^c\rangle,
\]
not a classical assignment of a permanent color to each named quark.}
\end{classroomqa}

\begin{classroomqa}
\audiencequestion{Will Bose condensation enter this course?
\lecturetimestamp{01:19:25}}
\lecturerresponse{Yes, through the Higgs phenomenon. A nonzero Higgs field
throughout the vacuum is analogous in important ways to a bosonic condensate.
The analogy will be developed after spontaneous symmetry breaking has been
introduced.\editorialnote{The Higgs vacuum expectation value is not literally an
ordinary gas of on-shell Higgs particles occupying one state. Condensation is
an illuminating analogy for a coherent bosonic field and a chosen vacuum,
not a complete microscopic identification.}}
\end{classroomqa}

\section{The Missing Neutral Weak Boson}

Return now to the group table. If the horizontal weak symmetry is \(SU(2)\),
then it has
\begin{equation}
  \dim SU(2)=3
\end{equation}
generators and therefore three gauge fields. The charged combinations
\(W^+\) and \(W^-\) account for only two. A third, neutral field is required.
It is closely related to the observed \(Z\) boson, but it is not simply the
photon. The precise neutral combinations involve both the third \(SU(2)_L\)
field and the \(U(1)_Y\) field. That mixing, together with the correct
definition of electric charge, is deliberately postponed to the next
lecture.

\section{Spontaneous Symmetry Breaking}

\subsection{A symmetric choice can have an unstable indecision}

Imagine a circular dinner table at which each place setting has a utensil to
the left and one to the right. Before anyone moves, left and right are
equivalent. Once one guest reaches to the right, every neighboring choice is
forced and the whole table settles into the right-handed arrangement. The
rules were symmetric, but the realized configuration selected one of two
possibilities. The story is only an analogy, yet it isolates the combination
we need: symmetry, multiple equivalent outcomes, and an instability that lets
an arbitrarily small disturbance select one outcome globally.

\subsection{The heads--tails lattice}

For a cleaner model, put a two-state variable \(s_i=\pm1\) at every site of a
lattice. Call \(+1\) heads and \(-1\) tails. Let neighboring sites prefer to
align:
\begin{equation}
  H=-J\sum_{\langle ij\rangle}s_i s_j,
  \qquad J>0.
\end{equation}
Parallel neighbors contribute lower energy than antiparallel neighbors.
There are therefore two degenerate ground states,
\begin{equation}
\begin{aligned}
  s_i&=+1\quad\text{for every }i,
  \\
  s_i&=-1\quad\text{for every }i.
\end{aligned}
\end{equation}
The Hamiltonian is invariant under the global flip \(s_i\to-s_i\), but either
chosen ground state is not.

Now fix one very distant boundary spin to heads. In an ordinary short-range
system we might expect one far-away degree of freedom to have no influence
here. In the ordered phase, however, that boundary condition selects the
all-heads vacuum. Once selected, a single tail in the sea of heads costs
energy because it disagrees with its neighbors. Local excitations are defined
relative to the chosen ground state even though the underlying equations
retain the heads--tails symmetry.

\subsection{Explicit breaking and domain walls}

Explicit breaking would add a local bias,
\begin{equation}
  H_{\rm explicit}
  =-J\sum_{\langle ij\rangle}s_i s_j-h\sum_i s_i,
  \qquad h\ne0.
\end{equation}
Then heads and tails do not have equal energy even at one isolated site, and
there is only one preferred vacuum. In spontaneous breaking \(h=0\): the
local rules are symmetric and the two vacua remain degenerate, but the system
occupies one of them.

Insist instead that the far-left boundary is all heads and the far-right
boundary all tails. The lowest-energy configuration must interpolate between
the two vacua. The interface is a domain wall.

\begin{figure}[htbp]
\centering
\includegraphics[width=0.9\textwidth]{lecture_06_figure_08.png}
\caption{The heads--tails lattice divided by a dashed domain wall. Opposite
boundary conditions force the system to contain an interface between the two
degenerate ordered states.}
\lectureframe{01:30:35}
\end{figure}

\begin{classroomqa}
\audiencequestion{What distinguishes spontaneous symmetry breaking from
ordinary explicit symmetry breaking?
\lecturetimestamp{01:28:36}}
\lecturerresponse{With explicit breaking the local Hamiltonian already assigns
different energies to heads and tails, so there is no pair of symmetry-related
ground states. With spontaneous breaking the equations retain the symmetry,
the vacua are degenerate, and boundary conditions or an arbitrarily small
disturbance select one. The possibility of a domain wall interpolating
between distinct vacua is a diagnostic of the latter situation.}
\end{classroomqa}

The model is deliberately primitive, but it contains the essential logic that
will be transferred to the Higgs field: a symmetric theory, an unstable or
degenerate vacuum structure, a selected ground state, and physical
consequences determined by that choice.

\section{Where the Lecture Leaves Us}

The fermion table turns the Standard Model group into a set of physical
directions. Color \(SU(3)\) moves vertically through red, green, and blue;
charged weak transformations move horizontally through quark and lepton
doublets; a phase symmetry supplies the preliminary \(U(1)\) factor. Gauge
boson emission realizes these group transformations as interactions.

The weakness of beta decay is then traced not to an exceptionally tiny
\(g_w\), but to the heavy mediator. In distance space, mass cuts off the
propagator beyond the Compton scale \(\hbar/(m_Wc)\). In momentum space,
\begin{equation}
  \mathcal A_{\rm weak}
  \sim\frac{g_w^2}{k^2+m_W^2}
  \simeq\frac{g_w^2}{m_W^2}
\end{equation}
at nuclear momenta. Old low-energy experiments measured this combination as
the Fermi constant.

The extended classroom discussion clarifies what this formula does and does
not mean. Internal \(W\) lines are virtual and off shell; energy--momentum is
still conserved. A hard enough probe can produce a real \(W\), but then it
has created a different final-state process. Collision products are not a
literal list of pre-existing constituents. Temperature changes the population
of energetic projectiles, photons can interact through charged virtual
particles, and Bose statistics is distinct from a dynamical force.

Finally, \(SU(2)\) requires a missing neutral gauge field related to the
\(Z\), and the heads--tails lattice supplies the first model of spontaneous
symmetry breaking. The next step is to combine those two unfinished threads:
electroweak mixing and the Higgs mechanism.
